\section{第一次访问匿名 VMA 时发生什么}

页表只能回答“处理器此刻怎样翻译这个地址”，不能回答“一个尚未驻留的地址
为什么有权在未来出现”。VMA 保存区间意图、来源和软件权限，PTE
保存当前驻留事实；再让 x86-64 页故障把真实用户访问转换成一次有边界、可
回滚的物理页提交。匿名映射、program break、受控用户栈和 Ring 3 用户堆由此
第一次形成完整生命周期。

\topic{上一阶段留下的虚拟内存缺口}
前面的 exec 已经能够从 rootfs 读取 ELF，建立候选地址空间，构造
\code{argc/argv/envp}，再由 PID1 组织 \code{spawn/exec/wait}。但每个合法
用户页都必须在进入 Ring 3 前存在：

\begin{itemize}
  \item ELF 的全部 \code{PT\_LOAD} 页在 exec 时立即分配；
  \item 用户栈固定映射 64 个 4 KiB 页；
  \item not-present 用户页故障统一表示进程错误；
  \item 用户程序没有匿名区间、program break 或对象级 heap。
\end{itemize}

这个模型可以验证隔离和映像事务，却把两个不同问题压进同一页表：

\begin{description}
  \item[地址意图] 进程是否拥有该范围，来源是什么，允许何种访问；
  \item[驻留事实] 此刻是否已有物理页，PTE 保存什么 frame 与硬件权限。
\end{description}

若没有第一层，处理器报告一个 not-present 地址时，Kernel 无法知道应该分配
零页、读取文件、复制 COW 页、增长栈，还是把它当作越界终止。

\topic{历史：从连续数据段到通用映射}
早期 Unix 进程常被画成 text、data、BSS、向高地址增长的 heap 和向低地址
增长的 stack。\term{program break}{进程数据段的逻辑末端}让用户分配器用
一个连续边界取得新空间：

\[
  B_{\mathrm{old}}\longrightarrow B_{\mathrm{new}},
  \qquad B_{\mathrm{base}}\le B_{\mathrm{new}}\le B_{\mathrm{limit}}.
\]

这个接口很适合一个主要动态区，却难以表达多个文件、共享对象、不同权限和
独立撤销。虚拟内存普及后，\code{mmap} 把“某个对象占据一段虚拟区间”变成
更一般的模型。匿名 mapping 没有文件来源；第一次访问时只需构造全零页。

把物理分配推迟到真实访问还有三个理由：

\begin{enumerate}
  \item 64 位虚拟空间远大于机器 RAM，预留不应立即等于物理承诺；
  \item 大栈通常只使用顶部少量页面；
  \item 文件页、COW 和零页都需要同一条 fault 解析入口。
\end{enumerate}

这里同时保留 brk 与匿名 mmap，先明确
连续区间和独立区间两种所有权。

\topic{x86-64 页故障提供的硬件事实}
用户 load、store 或 instruction fetch 首先查 TLB；未命中时由硬件遍历
PML4、PDPT、PD 和 PT。任何层级不存在、有效权限合取失败或保留位非法，
CPU 都会：

\begin{enumerate}
  \item 不提交本次内存访问；
  \item 把线性地址写入 CR2；
  \item 通过 IDT vector 14 进入异常入口；
  \item 压入返回现场并提供 page-fault error code。
\end{enumerate}

页故障处理使用的 error-code 位如下：

\begin{osgridlongtable}{L{0.08\textwidth}L{0.12\textwidth}L{0.31\textwidth}L{0.31\textwidth}}
位 & 名称 & 为零 & 为一\\ \hline
0 & P & 页不在场 & 已在场但权限违反\\ \hline
1 & W/R & 读访问 & 写访问\\ \hline
2 & U/S & supervisor 访问 & user 访问\\ \hline
3 & RSVD & 无保留位错误 & 页表保留位非法\\ \hline
4 & I/D & 数据访问 & 指令取址\\ \hline
\end{osgridlongtable}

RSVD 必须先于普通 not-present 策略拒绝，因为它通常表示 Kernel 构造了非法
页表。P 为一也不能被“补一张可写页”修复：向只读页写入正是硬件保护应捕获
的行为。

如果 Kernel 建好合法 PTE 并从异常返回，CPU 会重试原用户指令。用户程序不
需要显式循环，也不会从一个人为指定的“fault 后继”继续。

\topic{VMA 与 PTE 的正交关系}
考虑四页匿名区域，只触及第二页：

\begin{lstlisting}[caption={区间意图和驻留事实可以不同}]
VMA:
  [0x60000000, 0x60004000) Anonymous R|W

PTE:
  0x60000000 -> not present
  0x60001000 -> physical frame, P|U|W|NX
  0x60002000 -> not present
  0x60003000 -> not present
\end{lstlisting}

此时 virtual page 为四，resident page 只增加一。访问第三页时，VMA 允许
构造另一张零页；访问 \code{0x60005000} 时没有 VMA，即使它同样
not-present，也必须视为越界。

\begin{pdfdiagram}[node distance=8mm and 12mm]
  \node[os state] (vma) {VirtualMemoryMap\\区间、kind、R/W/X};
  \node[os hardware,right=of vma] (fault) {CR2 + \#PF error\\RIP/RSP};
  \node[os block,right=of fault] (policy) {fault policy\\允许或拒绝};
  \node[os state,below=of policy] (pte) {PTE\\frame + U/S/RW/NX};
  \node[os block,left=of pte] (frame) {Buddy frame\\direct-map 清零};
  \draw[os arrow] (vma) -- (policy);
  \draw[os arrow] (fault) -- (policy);
  \draw[os arrow] (policy) -- (pte);
  \draw[os arrow] (policy) -- (frame);
  \draw[os arrow] (frame) -- (pte);
\end{pdfdiagram}
\diagramnote{VMA 和硬件 fault 事实共同决定策略；页表与物理页是成功策略的
提交结果，不能反过来替代 VMA。}

\topic{用户虚拟地址窗口怎样划分}
当前使用以下固定窗口：

\begin{osgridlongtable}{L{0.25\textwidth}L{0.31\textwidth}L{0.32\textwidth}}
区域 & 范围或上限 & 语义\\ \hline
ELF 与 brk & 最高不越过 \code{0x60000000} & ELF eager，break 从最高段页尾开始\\ \hline
匿名窗口 & \code{[0x60000000,0x80000000)} & 512 MiB first-fit 或非覆盖 fixed\\ \hline
栈 guard & \code{0x00007FFFFF7EF000} & 永久无 VMA、无 PTE\\ \hline
用户栈 & \code{[0x00007FFFFF7F0000,0x00007FFFFFFF0000)} & 8 MiB 预留、向低地址提交\\ \hline
用户 heap & 最大 8 MiB & UserHeap 自行施加的运行时边界\\ \hline
\end{osgridlongtable}

这个匿名窗口不等于处理器全部用户地址能力。CPU profile 要求 48 位虚拟
地址；窄窗口是当前分配策略，使重叠、first-gap、容量耗尽和边界溢出
可以被完整测试。

\topic{描述符池为什么独立于每个 map}
\code{VirtualMemoryAreaPool} 拥有 8192 个实际描述符和一条自由索引链；
每个 \code{VirtualMemoryMap} 只链接带自己 owner identifier 的节点。单个
Process 最多拥有 4096 个区域。

这种设计有意避免：

\begin{itemize}
  \item 按 48 位地址空间的潜在页数建立位图；
  \item 在页故障路径递归调用尚可能依赖映射的动态分配器；
  \item 一个地址空间错误释放另一个地址空间的节点；
  \item 用“容量很大”掩盖元数据耗尽失败语义。
\end{itemize}

池持续记录：

\[
  \mathrm{free}+\mathrm{active}=\mathrm{capacity},
  \qquad
  \mathrm{release}\le\mathrm{acquire}.
\]

工作负载结束后还要求 active 为零、free 等于 capacity，并比较阶段内
acquire/release 增量。

\topic{Insert：先检查，再合并，最后申请}
VMA 使用页对齐半开区间 \([b,e)\)。Insert 先找到前驱和后继，再拒绝：

\[
  p.e>b\quad\text{或}\quad e>n.b.
\]

相邻区域只有 kind 与三项权限全部相等时才合并：

\begin{lstlisting}[caption={相邻区域的三种合并}]
[previous same] + [new] + [next same]
  -> extend previous through next, release next descriptor

[previous same] + [new]
  -> extend previous.end

[new] + [next same]
  -> lower next.begin
\end{lstlisting}

只有无法合并时才从池取得新节点。这样连续多次 brk growth 仍保持一段，而
不是每次系统调用永久消耗一个描述符。

\topic{Remove：中段拆分为何必须预取资源}
从 \([0,16)\) 移除 \([6,10)\) 会形成两个区域：

\[
  [0,6)\cup[10,16).
\]

区域数增加一。如果先把旧节点裁成左段，随后才发现没有新描述符，原图已经被
部分修改。因此 Remove 分两遍：

\begin{enumerate}
  \item 扫描全部重叠区间，确认 kind 一致并判断是否需要 split；
  \item 若需要 split，先 Acquire；
  \item 第二遍才执行完整删除、前缀/后缀裁剪或中段拆分。
\end{enumerate}

只有一个描述符的宿主测试刻意触发这一条件，要求返回
\code{AreaLimitExceeded} 且旧区域逐字段不变。

\topic{first-gap 是一项有界区间算法}
自动 mmap 从窗口起点开始，按页对齐计算候选地址。对每个有序 VMA：

\begin{enumerate}
  \item 已完全位于候选之前则跳过；
  \item 当前 VMA begin 与候选之间足够大则成功；
  \item 否则把候选推进到当前 VMA end 后再次对齐；
  \item 任何加法、对齐或 \code{end-candidate} 检查失败都明确返回。
\end{enumerate}

当前实现为线性复杂度。这是经过选择的学习基线：公开接口没有泄漏链表，
未来可以用树替换，而本阶段所有修改路径仍能在源码中直接推演。

\topic{系统调用 39 至 42}
共享 ABI 在旧编号之后追加：

\begin{osgridlongtable}{L{0.08\textwidth}L{0.29\textwidth}L{0.31\textwidth}L{0.20\textwidth}}
号 & 名称 & 输入 & 成功结果\\ \hline
39 & \code{MapAnonymousMemory} & 地址、长度、权限、flags & 映射起点\\ \hline
40 & \code{UnmapMemory} & 页对齐地址、长度 & 0\\ \hline
41 & \code{SetProgramBreak} & 0 查询或新地址 & 当前 break\\ \hline
42 & \shortstack[l]{\code{GetVirtualMemory}\\\code{Statistics}} &
目标地址、结构尺寸 & 0\\ \hline
\end{osgridlongtable}

AMD64 用户包装继续使用：

\begin{lstlisting}[caption={四参数系统调用寄存器}]
RAX = system call number
RDI = argument 0
RSI = argument 1
RDX = argument 2
R10 = argument 3
\end{lstlisting}

fixed mapping 不是破坏性 Linux \code{MAP\_FIXED}：它只在指定页对齐空洞
成功，重叠返回 \code{AddressInUse}。当前不承担跨 kind 替换事务。

权限只接受 NONE、R、R|W 与 R|X。W 或 X 缺少 R、W|X 和未知位都拒绝；
VMA 软件权限随后成为 PTE RW/NX 的唯一策略来源。

\topic{统一缺页提交路径}
\code{HandleUserPageFault} 的判断顺序是安全边界本身：

\begin{lstlisting}[caption={用户页故障决策顺序}]
require U/S == user
reject RSVD
reject P == present violation
decode write and instruction access
find VMA containing CR2 page
check readable/writable/executable
check stack policy when kind == UserStack
reject missing eager ExecutableImage
allocate and map one page
zero the complete frame through direct map
update counters
return to retry the user instruction
\end{lstlisting}

\begin{failurebox}
把 P=1 的保护 fault 当作普通缺页，会让只读映射在首次写时被静默升级；把
没有 VMA 的 not-present fault 当作匿名页，会让任意越界地址变成合法内存；
把 RSVD 当作用户错误，又会掩盖 Kernel 页表损坏。三者必须在分配 frame
之前分开。
\end{failurebox}

\topic{一页提交为何必须可回滚}
\code{MapDemandPage} 取得 frame、建立必要页表分支、写 PTE，再通过高半区
direct-map 清零 4096 字节。若 direct-map 地址无效或提交失败，必须调用
\code{ReleaseUserPage}：

\begin{itemize}
  \item 清除刚写入的用户 PTE；
  \item 释放数据 frame；
  \item 释放本事务造成的空 PT/PD/私有 PDPT；
  \item 不增加 resident 与成功 fault 统计。
\end{itemize}

用户指令尚未从异常返回，因此不会观察清零前内容。未来 SMP 会额外需要并发
fault 协调与远端 TLB shootdown；当前单 BSP 不把这项未来承诺伪装成已完成。

\topic{栈预留和实际提交必须分开}
装载 ELF 时，完整 8 MiB \code{UserStack} VMA 已经存在；参数布局只提交
实际覆盖的顶部页。设栈顶为 \(T\)，最低已提交地址为 \(C\)，fault page 为
\(F\)，一次合法增长首先要求：

\[
  F+4096=C.
\]

此外还要求：

\[
  F\ge B_{\mathrm{stack}},
  \qquad
  A_{\mathrm{fault}}\le RSP+4096,
  \qquad
  RSP\le A_{\mathrm{fault}}+65536.
\]

\begin{pdfdiagram}[node distance=5mm]
  \node[os hardware] (top) {stack top};
  \node[os block,below=of top] (resident) {已提交参数与调用栈页};
  \node[os state,below=of resident] (candidate) {唯一合法下一页\\fault page};
  \node[os block,below=of candidate] (reserve) {其余 stack VMA\\仍 not-present};
  \node[os hardware,below=of reserve] (guard) {永久 guard\\无 VMA};
  \draw[os arrow] (top) -- (resident);
  \draw[os arrow] (resident) -- (candidate);
  \draw[os arrow] (candidate) -- (reserve);
  \draw[os arrow] (reserve) -- (guard);
\end{pdfdiagram}
\diagramnote{栈 VMA 内也不能任意跳跃建页；连续 committed bottom 和保存的
用户 RSP 共同提供合法增长证据。}

guard 完全没有 VMA，而不只是一个临时 not-present PTE。因此通用
\code{FindContaining} 已经能够拒绝它，后续文件页或 COW 策略也不会误认。

\topic{Kernel 用户复制为什么只解析部分 VMA}
系统调用经常需要读用户请求或写用户结果。如果目标是合法、尚未驻留的匿名/
break 页，直接返回 bad pointer 会让按需语义在 syscall 内外不一致。
\code{CopyFromUser} 与 \code{CopyToUser} 因此可以调用
\code{ResolveNonStackPage}。

它们不能增长用户栈。Kernel copy 没有一条正在执行的用户 push/store，也没有
可用于证明栈推进的异常现场；替用户制造栈页会绕过 RSP 策略。候选 exec 参数
栈则由已经验证最低地址的 \code{PrepareUserStack} 显式预提交。

\topic{brk 只管理页级承诺}
ELF 最高 \code{PT\_LOAD} 末端向上页对齐得到 break base。查询使用请求零；
扩大时插入或合并 \code{ProgramBreak} VMA，不分配数据页。缩小时按页对齐
边界撤销 VMA 并释放实际驻留页。

用户程序可能请求非页对齐 break。逻辑 break 精确保留请求值，VMA 则覆盖
到向上对齐的页尾：

\[
  E_{\mathrm{VMA}}=\operatorname{alignUp}(B,4096).
\]

这让对象分配器使用字节边界，同时让页表仍只处理整页。

\topic{unmap 如何同时回收数据页和页表页}
\code{UnmapMemory} 只接受 Anonymous kind，不能顺便删除 ELF、stack 或
program break。完成 VMA remove 后逐页查询：

\begin{description}
  \item[NotMapped] 该页只是 reservation，不存在 frame，直接继续；
  \item[Mapped] 释放用户 frame，清 PTE，并向上回收空页表分支；
  \item[其他状态] 地址空间内部不一致，不能伪装成成功。
\end{description}

相邻页共享 PT/PD，因此“撤销一页”不对应固定数量页表 frame。实现累计
\code{ReleaseUserPage} 返回的真实 reclaimed count，而不是用公式猜测。

exec 成功后销毁旧页表和旧 VMA；Process 退出切回永久 Kernel CR3 后销毁
当前地址空间。最终 VMA 池与 frame、buddy、KVA、KernelStack、fd 和 VFS
资源一起比较基线。

\topic{UserHeap 为什么放在 Ring 3}
Kernel 提供页级 break，不应理解每个 C++ 对象的大小和释放顺序。用户
\code{UserHeap} 注入 break operation、页大小、增长 quantum 与最大容量，
再在取得的连续空间中管理块。

每个 64 字节 header 显式保存：

\begin{itemize}
  \item payload capacity 与用户 requested size；
  \item 前一物理块 offset；
  \item 前后空闲块 offset；
  \item allocated/free 状态、signature 与 reserved。
\end{itemize}

块按 16 字节对齐，关系使用相对 heap base 的 \code{uint64\_t} offset，
不依赖宿主指针宽度容器。

\begin{pdfdiagram}[node distance=4mm and 5mm]
  \node[os state] (h0) {Header A\\allocated};
  \node[os block,right=of h0] (p0) {payload A};
  \node[os state,right=of p0] (h1) {Header B\\free};
  \node[os block,right=of h1] (p1) {free payload};
  \node[os state,below=of h1] (free) {free-list prev/next};
  \draw[os arrow] (h0) -- (p0);
  \draw[os arrow] (p0) -- (h1);
  \draw[os arrow] (h1) -- (p1);
  \draw[os arrow] (h1) -- (free);
\end{pdfdiagram}
\diagramnote{物理相邻关系用于 coalesce，空闲双链用于 first-fit；同一个
free block 同时属于两张图。}

\topic{first-fit、split 与双向 coalesce}
Allocate 把请求向上对齐到 16 字节，沿 free list 找第一个足够块。若没有，
按页对齐 growth quantum 扩大 break，再重新查找。

只有余量满足：

\[
  R\ge \operatorname{sizeof(Header)}+16
\]

才 split。否则整个容量交给本次 allocation，但统计仍保存原始 requested
size，不能把内部碎片算成用户请求。

Release 只接受精确 payload 起点。标记 free 后先合并前一物理邻居，再合并
后一邻居；将消失的节点必须先从 free list 摘除，最终合并块只在链中出现一次。
外部指针、块内部指针、重复释放和 signature 错误分别返回明确状态。

\topic{Validate 为什么要遍历两张图}
只验证 free list 不能发现一个活动块覆盖下一 header；只验证物理块又不能
发现 free 节点重复或遗漏。Validate 首先从 heap base 走到 capacity：

\begin{itemize}
  \item header 完整落在范围内；
  \item signature、state、16 字节容量对齐有效；
  \item previous physical offset 与观察值一致；
  \item allocated 请求非零且不超过容量；
  \item next offset 严格前进且最终恰好等于 capacity。
\end{itemize}

随后遍历 free list，核对 prev/next、状态、节点数和物理遍历观察到的 free
数完全相同。活动 allocation 数和 requested bytes 也从事实重建后比较。

\topic{日志为什么采用二次幂采样}
一个 32 MiB 区间有 8192 页。若每次 fault 都打印 CR2、frame 和页表层级，
115200 波特串口会成为主要工作负载，并改变 PIT 抢占与 TCG 时序。

Kernel 只在每 Process 累计值达到二次幂时输出：

\begin{lstlisting}[caption={按需分页采样日志}]
[OS][KERNEL][VM] DEMAND_FAULT_COUNT=0x...
[OS][KERNEL][VM] STACK_GROWTH_COUNT=0x...
\end{lstlisting}

不同 Process 的计数各自从一开始，整机可以出现多行同值。宿主 runner 因此
检查“至少出现且数值非零”，不要求精确 multiplicity。最终再一次输出池容量、
active、peak、acquire 和 release。

\topic{分别在 VMA、PTE、随机模型和 QEMU 中检查什么}\par
\begin{osgridlongtable}{L{0.17\textwidth}L{0.71\textwidth}}
层次 & 直接证明\\ \hline
VMA 单元 & 排序、重叠、三类合并、trim/split、kind guard、first-gap、耗尽失败原子性\\ \hline
Heap 单元 & 增长、对齐、split、复用、双向合并、耗尽、外部和重复释放\\ \hline
固定随机 & VMA 与 heap 各 100000 步，逐步对照独立参考模型与活动载荷\\ \hline
组合生命周期 & 128 轮 reservation、真实 PTE、空页表回收和 frame/descriptor 基线\\ \hline
Ring 3 & 32 MiB 稀疏触页、2 MiB brk、递归栈和 5000 步 UserHeap\\ \hline
保护与容量 & guard/只读写 vector 14；64 MiB、256 MiB、64 GiB 同一工作负载\\ \hline
\end{osgridlongtable}

reservation 测试要求 VMA 插入后 frame count 不变；只看到 mmap 返回成功并
不能证明这一点。QEMU 又证明真实 CR2、异常入口、CR3、PTE 与用户指令重试；
宿主模型不能替代这层硬件证据。

\topic{代码阅读顺序}\par
\begin{enumerate}
  \item \filepath{source/abi/include/os/abi/virtual\_memory.hpp}；
  \item \filepath{memory/virtual\_memory\_area.hpp/.cpp}；
  \item \filepath{user/user\_memory.hpp/.cpp}；
  \item \filepath{arch/exceptions.cpp} 到 ProcessRuntime fault 入口；
  \item \filepath{user/system\_calls.cpp} 与 Ring 3 包装；
  \item \filepath{source/user/include/os/user/user\_heap.hpp} 及实现；
  \item VMA、heap 单元和随机测试；
  \item user VM 生命周期集成测试；
  \item memory、guard 与 protection 三个用户 probe；
  \item QEMU 最终 VMA 与 ProcessTree 资源摘要。
\end{enumerate}

调试时在 \code{HandleUserPageFault} 断住，记录 CR2、error code、保存的
RIP/RSP、命中 VMA、分配 frame 和 resident 前后值。继续执行后原 RIP 应
重试成功，而不是跳到人工恢复点。

\topic{阶段边界为何仍然重要}
当前没有实现：

\begin{itemize}
  \item file-backed VMA、按需 ELF 或 clean page cache；
  \item \code{MAP\_PRIVATE}、\code{MAP\_SHARED} 或 \code{mprotect}；
  \item fork、COW、匿名页引用计数、swap 或 overcommit；
  \item ASLR、动态链接器、共享库或完整 libc；
  \item 多 Thread heap arena、size class 或尾块自动缩 break。
\end{itemize}

文件页需要 \code{(Vnode,page index)} 身份、尾页清零、缓存共享和
truncate/write 失效；COW 又需要 VMA 权限、PTE COW 状态和物理页引用共同
成立。因此下一章先在 fault 边界上加入文件来源，再加入 fork/COW。

\topic{本阶段完成判定}
一次用户页故障后继续运行只覆盖了正常路径。测试还要同时满足：

\begin{enumerate}
  \item VMA 图与描述符池在所有操作后通过 Validate；
  \item 未触及 reservation 不消耗数据 frame；
  \item 首次匿名读为零，写入在生命周期内保持；
  \item stack 只连续增长，guard 与只读写分别被 vector 14 捕获；
  \item unmap、break shrink、exec 和 exit 回收驻留页与空页表；
  \item UserHeap 随机结束后活动 allocation 为零；
  \item 十一个 Process 生命周期后 VMA active 为零，acquire/release 守恒；
  \item 64 MiB、256 MiB、64 GiB 三档使用同一实现通过；
  \item 单元、集成、随机、产物、故障和历史回归同时保持；
  \item 教材中的容量、日志和状态转换与精确源码及运行结果一致。
\end{enumerate}

这些条件把“能跑”提升为一套可以为文件页和 COW 继续依赖的虚拟内存契约。
